% GENERATED FILE. Edit canonical lessons, then run npm run generate:books.
% canonical-source-hash: fnv1a64:15a39b0a8326c96e
% canonical-lessons: FA-C387-mohajerat, FA-C387-mohajer, FA-C387-aqalliyyat, FA-C387-aksariyyat, FA-C387-tabiz

\chapter{Migration and minorities}
\label{ch:fa-migration-and-minorities}
\hlchaptermodality{387}

\begin{chapteropening}
\textbf{By the end of this chapter:} \emph{I can talk about migration and migrants, minorities and majorities, and discrimination.}

Complete the last lesson of chapter 387: I can talk about migration and migrants, minorities and majorities, and discrimination.
\end{chapteropening}

\section[\texorpdfstring{\fa{مهاجرت} (mohâjerat)}{مهاجرت (mohâjerat)}]{\fa{مهاجرت} (mohâjerat) — migration, emigration}

\label{lesson:FA-C387-mohajerat}



\begin{quote}
Before the new one: say the Persian for poverty, then the Persian for wealth.
\end{quote}

\subsection*{You'll want to know: \fa{مهاجرت}}
\textbf{\fa{مهاجرت}} (\emph{mohâjerat}) — \textquotedblleft{}migration, emigration\textquotedblright{}.

\begin{quote}
\fa{مهاجرت} \fa{به} \fa{خارج} \fa{برای} \fa{او} \fa{آسان} \fa{نبود}.
\end{quote}

\begin{quote}
\emph{mohâjerat be khârej barâye u âsân nabud.}
\end{quote}

\begin{quote}
Moving abroad was not easy for her.
\end{quote}

\emph{Mohâjerat kardan} — to emigrate.

\begin{cousinweb}
\textbf{\fa{مهاجرت}} is usually traced to Arabic \emph{muhājara} \textquotedblleft{}emigration\textquotedblright{}, from the same Arabic family as \emph{hijra} \textquotedblleft{}departure\textquotedblright{}, which EtymDB lists.
\end{cousinweb}

\begin{grammarlens}[title={What you've built}]
One more word for telling, explaining and describing at B1.
\end{grammarlens}

\subsection*{Your turn}
\begin{itemize}
\raggedright
  \item \emph{Say it:} \emph{mohâjerat}
  \item \emph{Say it:} \emph{mohâjerat}, once more
  \item \emph{Say it:} a sentence of your own with \emph{mohâjerat}
  \item \emph{Recall:} say the Persian for poverty, then the Persian for wealth, then say \emph{mohâjerat} again
\end{itemize}

\begin{tcolorbox}[breakable,colback=teal!4,colframe=teal!35!black,title={Before you move on}]
What does \fa{مهاجرت} mean? (\textquotedblleft{}Migration, emigration\textquotedblright{}.) Say it once more.
\end{tcolorbox}
\section[\texorpdfstring{\fa{مهاجر} (mohâjer)}{مهاجر (mohâjer)}]{\fa{مهاجر} (mohâjer) — a migrant, an immigrant}

\label{lesson:FA-C387-mohajer}



\begin{quote}
Before the new one: say the Persian for wealth, then the Persian for migration.
\end{quote}

\subsection*{You'll want to know: \fa{مهاجر}}
\textbf{\fa{مهاجر}} (\emph{mohâjer}) — \textquotedblleft{}a migrant, an immigrant\textquotedblright{}.

\begin{quote}
\fa{خیلی} \fa{از} \fa{مهاجرها} \fa{در} \fa{این} \fa{محله} \fa{زندگی} \fa{میکنند}.
\end{quote}

\begin{quote}
\emph{kheyli az mohâjerhâ dar in mahalle zendegi mikonand.}
\end{quote}

\begin{quote}
Many migrants live in this neighbourhood.
\end{quote}

Plural \emph{mohâjerân} or \emph{mohâjerhâ}.

\begin{cousinweb}
\textbf{\fa{مهاجر}} comes from Arabic \emph{muhājir} \textquotedblleft{}emigrant\textquotedblright{} (EtymDB): one who has left home.
\end{cousinweb}

\begin{grammarlens}[title={What you've built}]
One more word for telling, explaining and describing at B1.
\end{grammarlens}

\subsection*{Your turn}
\begin{itemize}
\raggedright
  \item \emph{Say it:} \emph{mohâjer}
  \item \emph{Say it:} \emph{mohâjer}, once more
  \item \emph{Say it:} a sentence of your own with \emph{mohâjer}
  \item \emph{Recall:} say the Persian for wealth, then the Persian for migration, then say \emph{mohâjer} again
\end{itemize}

\begin{tcolorbox}[breakable,colback=teal!4,colframe=teal!35!black,title={Before you move on}]
What does \fa{مهاجر} mean? (\textquotedblleft{}A migrant, an immigrant\textquotedblright{}.) Say it once more.
\end{tcolorbox}
\section[\texorpdfstring{\fa{اقلیت} (aqalliyat)}{اقلیت (aqalliyat)}]{\fa{اقلیت} (aqalliyat) — a minority}

\label{lesson:FA-C387-aqalliyyat}



\begin{quote}
Before the new one: say the Persian for migration, then the Persian for a migrant.
\end{quote}

\subsection*{You'll want to know: \fa{اقلیت}}
\textbf{\fa{اقلیت}} (\emph{aqalliyat}) — \textquotedblleft{}a minority\textquotedblright{}.

\begin{quote}
\fa{آنها} \fa{در} \fa{این} \fa{کشور} \fa{یک} \fa{اقلیت} \fa{کوچک} \fa{هستند}.
\end{quote}

\begin{quote}
\emph{ânhâ dar in keshvar yek aqalliyat-e kuchek hastand.}
\end{quote}

\begin{quote}
They are a small minority in this country.
\end{quote}

Plural \emph{aqalliyathâ}. The opposite is \emph{aksariyat}.

\begin{cousinweb}
\textbf{\fa{اقلیت}} is usually traced to Arabic \emph{aqalliyya} \textquotedblleft{}minority\textquotedblright{}, built on \emph{aqall} \textquotedblleft{}fewer, least\textquotedblright{}.
\end{cousinweb}

\begin{grammarlens}[title={What you've built}]
One more word for telling, explaining and describing at B1.
\end{grammarlens}

\subsection*{Your turn}
\begin{itemize}
\raggedright
  \item \emph{Say it:} \emph{aqalliyat}
  \item \emph{Say it:} \emph{aqalliyat}, once more
  \item \emph{Say it:} a sentence of your own with \emph{aqalliyat}
  \item \emph{Recall:} say the Persian for migration, then the Persian for a migrant, then say \emph{aqalliyat} again
\end{itemize}

\begin{tcolorbox}[breakable,colback=teal!4,colframe=teal!35!black,title={Before you move on}]
What does \fa{اقلیت} mean? (\textquotedblleft{}A minority\textquotedblright{}.) Say it once more.
\end{tcolorbox}
\section[\texorpdfstring{\fa{اکثریت} (aksariyat)}{اکثریت (aksariyat)}]{\fa{اکثریت} (aksariyat) — a majority}

\label{lesson:FA-C387-aksariyyat}



\begin{quote}
Before the new one: say the Persian for a migrant, then the Persian for a minority.
\end{quote}

\subsection*{You'll want to know: \fa{اکثریت}}
\textbf{\fa{اکثریت}} (\emph{aksariyat}) — \textquotedblleft{}a majority\textquotedblright{}.

\begin{quote}
\fa{اکثریت} \fa{مردم} \fa{با} \fa{این} \fa{طرح} \fa{موافق} \fa{هستند}.
\end{quote}

\begin{quote}
\emph{aksariyat-e mardom bâ in tarh movâfeq hastand.}
\end{quote}

\begin{quote}
The majority of people agree with this plan.
\end{quote}

The opposite is \emph{aqalliyat}.

\begin{cousinweb}
\textbf{\fa{اکثریت}} is usually traced to Arabic \emph{akthariyya} \textquotedblleft{}majority\textquotedblright{}, built on \emph{akthar} \textquotedblleft{}more, most\textquotedblright{}; Persian says the Arabic \emph{th} as \emph{s}.
\end{cousinweb}

\begin{grammarlens}[title={What you've built}]
One more word for telling, explaining and describing at B1.
\end{grammarlens}

\subsection*{Your turn}
\begin{itemize}
\raggedright
  \item \emph{Say it:} \emph{aksariyat}
  \item \emph{Say it:} \emph{aksariyat}, once more
  \item \emph{Say it:} a sentence of your own with \emph{aksariyat}
  \item \emph{Recall:} say the Persian for a migrant, then the Persian for a minority, then say \emph{aksariyat} again
\end{itemize}

\begin{tcolorbox}[breakable,colback=teal!4,colframe=teal!35!black,title={Before you move on}]
What does \fa{اکثریت} mean? (\textquotedblleft{}A majority\textquotedblright{}.) Say it once more.
\end{tcolorbox}
\section[\texorpdfstring{\fa{تبعیض} (tab'iz)}{تبعیض (tab'iz)}]{\fa{تبعیض} (tab'iz) — discrimination}

\label{lesson:FA-C387-tabiz}



\begin{quote}
Before the new one: say the Persian for a minority, then the Persian for a majority.
\end{quote}

\subsection*{You'll want to know: \fa{تبعیض}}
\textbf{\fa{تبعیض}} (\emph{tab'iz}) — \textquotedblleft{}discrimination\textquotedblright{}.

\begin{quote}
\fa{تبعیض} \fa{بین} \fa{دختر} \fa{و} \fa{پسر} \fa{درست} \fa{نیست}.
\end{quote}

\begin{quote}
\emph{tab'iz beyn-e dokhtar o pesar dorost nist.}
\end{quote}

\begin{quote}
Discrimination between girls and boys is wrong.
\end{quote}

Formal. \emph{Tab'iz-e nezhâdi} is racial discrimination.

\begin{cousinweb}
\textbf{\fa{تبعیض}} is usually traced to Arabic \emph{tab'īḍ} \textquotedblleft{}dividing into parts\textquotedblright{}: discrimination splits people into groups and treats them apart.
\end{cousinweb}

\begin{grammarlens}[title={What you've built}]
One more word for telling, explaining and describing at B1.
\end{grammarlens}

\subsection*{Your turn}
\begin{itemize}
\raggedright
  \item \emph{Say it:} \emph{tab'iz}
  \item \emph{Say it:} \emph{tab'iz}, once more
  \item \emph{Say it:} a sentence of your own with \emph{tab'iz}
  \item \emph{Recall:} say the Persian for a minority, then the Persian for a majority, then say \emph{tab'iz} again
\end{itemize}

\begin{tcolorbox}[breakable,colback=teal!4,colframe=teal!35!black,title={Before you move on}]
What does \fa{تبعیض} mean? (\textquotedblleft{}Discrimination\textquotedblright{}.) Say it once more.
\end{tcolorbox}
